\chapter{Introduction}
\label{Chapter1}

\section{Motivation}
\label{sec:ch1-motivation}

Large language models (LLMs) are neural networks that are trained to predict the next token of text and that, at sufficient scale, can perform translation, question answering, summarisation and many other language tasks from a textual prompt alone, without task-specific training. Nearly all current LLMs are built on the transformer, an architecture that replaced recurrence with an attention mechanism and that was introduced for machine translation. The capability of these models has grown with their size: openly released model families such as Llama~2 span 7 to 70 billion parameters.

The same scale that makes these models capable makes them expensive to run. Every parameter must be stored in memory and read for every generated token, so the size of a model sets both its memory footprint and its generation speed. Lin et al.\ describe the problem directly: the model size raises the hardware barrier for serving, because the weights must fit in memory, and slows down token generation, because generation is limited by memory bandwidth \cite{lin2024awq}. In generative inference the output is produced one token at a time, and at small batch sizes and sequence lengths the time needed to load the weights from memory dominates the time of each step. On top of the weights, the model keeps a key--value (KV) cache of intermediate attention state for every previous token; this cache is stored in GPU memory in addition to the parameters and grows linearly with the sequence length and the batch size \cite{zhang2023h2o}, and it is large enough that the way a serving system manages it limits how many requests can be batched together.

These costs matter most where resources are limited: on edge and embedded devices, on single consumer GPUs, and in services that must answer within a fixed time. In such settings a model is only useful if it fits in the available memory and answers quickly enough, so there is strong practical interest in making LLMs smaller and faster without losing much of their quality. Model compression is the family of techniques that pursues this goal. A recent survey groups them into quantization, which stores weights and activations with fewer bits; pruning, which removes parameters or structures; and knowledge distillation, which trains a smaller model to imitate a larger one. Post-training methods are especially attractive because they do not require retraining the full model. LLM.int8() showed that the feed-forward and attention projections of models with up to 175 billion parameters can be run with 8-bit matrix multiplication, halving inference memory, once the few emergent outlier features are handled separately in 16 bits \cite{dettmers2022llmint8}; GPTQ (published as OPTQ) quantizes weights to 3 or 4 bits using approximate second-order information and can process a 175-billion-parameter model in about four GPU hours \cite{frantar2023gptq}; AWQ protects the one percent of weights that matter most, as identified from the activations, and keeps low-bit weight-only quantization accurate without back-propagation \cite{lin2024awq}; SmoothQuant moves the difficulty of quantizing activation outliers into the weights so that both weights and activations can be stored in 8 bits \cite{xiao2023smoothquant}; Wanda prunes a model with a simple weight-times-activation criterion and no weight update \cite{sun2024wanda}; and H2O shrinks the KV cache by keeping only the recent tokens and the few tokens that receive most of the attention \cite{zhang2023h2o}.

Each of these methods is effective on its own terms. What is less clear is how they compare with one another, how much of their reported benefit survives in the inference engine that will actually serve the model, and whether a model should be compressed uniformly or differently in different parts. This thesis is motivated by those questions.

\section{Problem Statement}
\label{sec:ch1-problem}

Compression techniques for LLMs are usually published as individual methods, each with its own model, calibration data, evaluation data, hardware and software. Their results are therefore difficult to compare, and a practitioner who must choose a technique and a setting for a particular model and device has little controlled evidence to rely on. The survey that preceded this thesis, written by the author and the supervisor and summarised in Section~\ref{sec:ch1-scope}, reviewed about fifty papers and identified five gaps that the present work takes as its starting point. First, there is no practical comparison of quantization precision levels under identical conditions. Second, there are no clear rules for choosing between quantization-only, pruning-only and combined pipelines. Third, pruning reduces the number of parameters but does not always reduce real latency. Fourth, there is little deployment-aware evaluation of the same compressed model across different inference tools such as Hugging Face Transformers, llama.cpp, vLLM and TensorRT-LLM. Fifth, there is little adaptive compression that chooses the bit width, pruning ratio or compression level according to a latency, memory or accuracy requirement.

A further, more specific problem underlies the fifth gap. Almost all post-training methods apply one setting, such as 4-bit weights or 50 percent sparsity, to every layer of the model. Yet the layers of a transformer are not equally important, and a method that could measure how much each layer tolerates could compress the tolerant layers more and the sensitive layers less. Whether such a per-layer plan actually improves the trade-off between quality, size and speed, and whether a cheap sensitivity measurement is a good enough guide, has not been established under controlled conditions.

The problem addressed by this thesis is therefore to build a single framework in which published compression techniques for weights, activations and the KV cache can be run with and without a sensitivity-driven per-layer plan, under one protocol, and in which every resulting model is measured not only in a simulation of the compressed arithmetic but also in the inference engines that would deploy it, so that the five gaps above can be examined with directly comparable evidence.

\section{Aims and Objectives}
\label{sec:ch1-aims}

The aim of this thesis is to determine whether sensitivity-driven, per-layer compression planning improves the quality, size and latency trade-off of post-training compression for a decoder-only LLM, and to characterise how compression choices translate into real deployment metrics across inference engines. The aim is pursued through the following research questions.

\begin{enumerate}
  \item \textbf{RQ1.} How unevenly is sensitivity to compression distributed across the layers of a decoder-only LLM, and does a cheap measured score, namely the rise in perplexity when a layer is bypassed, rounded to fewer bits, or has its KV cache reduced, predict the damage done by compressing that layer?
  \item \textbf{RQ2.} Under identical calibration data, seed, evaluation data, hardware and software, does a sensitivity-driven per-layer plan that assigns 8 bits to sensitive layers and 4 bits to the rest improve the quality, size and latency trade-off of published weight, activation and KV-cache compression techniques compared with their usual uniform settings?
  \item \textbf{RQ3.} When is quantization-only, pruning-only or a combined pipeline preferable, given that pruned weights are stored as zeros and therefore bring no size or latency gain in the inference engines studied here?
  \item \textbf{RQ4.} How do simulated compression results obtained in Hugging Face Transformers carry over to real file size, memory and decode latency in llama.cpp, vLLM and TensorRT-LLM?
  \item \textbf{RQ5.} Can a constrained multi-objective search over techniques and their settings find configurations that satisfy a deployment requirement on perplexity and decode latency while minimising perplexity, latency and model size?
\end{enumerate}

To answer these questions the following objectives were set.

\begin{enumerate}
  \item \textbf{O1.} Build a per-layer sensitivity profiler and planner that measures the effect of bypassing, rounding and KV-cache reduction for each decoder layer, and that produces weight, activation, joint weight--activation, pruning and KV-cache plans together with a prediction of the memory each plan would need.
  \item \textbf{O2.} Implement or adapt published quantization, pruning, low-rank and KV-cache techniques so that each can be run with and without the plan on identical code.
  \item \textbf{O3.} Define one comparison protocol, consisting of the FP16 model, each technique at its usual uniform setting, and each technique following the plan, with the same calibration data, seed, evaluation data, hardware and software throughout.
  \item \textbf{O4.} Export every compressed model to llama.cpp, vLLM and TensorRT-LLM and measure its real file size, memory use and latency in each engine.
  \item \textbf{O5.} Implement a constrained multi-objective search that selects the technique and settings for each compression stage under a deployment requirement, and report its outcomes.
  \item \textbf{O6.} Analyse the results against the five gaps and derive guidance on precision levels, pipeline choice, pruning, engine behaviour and adaptive planning for practitioners.
\end{enumerate}

\section{Scope and Limitations}
\label{sec:ch1-scope}

The work builds on the survey by the author and the supervisor, which reviewed around fifty papers published from 2020 onwards in five groups (quantization, pruning and sparsity, hybrid quantization and pruning, knowledge distillation, and KV-cache compression) and which proposed a sensitivity-driven, adaptive framework. The thesis implements and evaluates that framework. The scope is deliberately bounded as follows.

The thesis considers post-training compression only. No full retraining, quantization-aware training or knowledge distillation is performed, and no custom hardware or new GPU kernels are developed. The techniques studied are published methods, most of them re-implemented in simplified form within the framework so that they can follow a per-layer plan; the aim is a controlled comparison of the planning idea, not a reproduction of each method's best published number. These simplified versions were not compared against the original implementations, so their results may differ from the published ones.

A single model is used throughout: TinyLlama, a 1.1-billion-parameter model that follows the architecture and tokenizer of Llama~2. It was chosen because it is small enough to be compressed, exported and measured many hundreds of times on the available hardware while still being a genuine decoder-only LLM of the Llama family. Findings may not transfer unchanged to other models, in particular to much larger ones.

Each configuration is run with one seed, so the variation between runs with different seeds is not measured. Quality is measured by perplexity on WikiText-2 and, for selected final models, by zero-shot accuracy on four common-sense benchmarks. Deployment metrics are measured in four engines: Hugging Face Transformers, which evaluates the quality of the simulated compression, and llama.cpp, vLLM and TensorRT-LLM, which give real file size, memory and speed after re-rounding the weights onto their own grids. Perplexity measured in an engine can differ from perplexity measured in Hugging Face (for one compressed model in this work, 11.50 in Hugging Face and 13.0 in vLLM), so the search judges perplexity in vLLM. The hardware is an NVIDIA A100 80~GB GPU on Google Colab and a desktop PC with an NVIDIA RTX~5070~Ti 16~GB GPU. Speed is measured on these two GPUs only, and the ranking of methods by speed may differ on other hardware.

Several techniques were attempted and dropped because they could not be made to run reliably in the framework, namely AQLM, QTIP and ABQ-LLM. BiLLM is included without a planned variant, and the planned variant of RPTQ was dropped. Pruned weights are stored as zeros rather than in a sparse format, so the pruning experiments measure quality only; they cannot show size or latency gains in the engines, and the thesis does not claim any. Finally, the multi-objective search is reported as a set of outcomes rather than as a pass or fail, since there is no single correct answer to a constrained trade-off.

\section{Structure of the Thesis}
\label{sec:ch1-structure}

The remainder of the thesis is organised as follows. Chapter~\ref{Chapter2} gives the background needed to read the later chapters: what a language model is and how it is measured, the transformer and the decoder-only architecture, how inference works and why decoding is limited by memory bandwidth, the model families and in particular TinyLlama, the main compression techniques, the inference engines, the evaluation metrics, and the basics of multi-objective optimisation. Chapter~\ref{Chapter3} reviews the published literature in groups, compares the groups, and states the research gaps that this thesis responds to. Chapter~4 describes the framework and the experimental method, Chapter~5 reports and discusses the results, and Chapter~6 concludes and suggests future work.
